\chapter*{TD1 : Variables aléatoires réelles et fonctions de répartition}

\setcounter{exercice}{0}

\begin{mdframed}
\begin{exercice}
Soit $X$ une variable aléatoire dont la fonction de répartition $F_X$ est définie par :

$$\forall t \in \mathbb{R} \quad F_X(t) = \begin{cases} 
0 & \text{si } t < -2, \\ 
\frac{1}{30} & \text{si } -2 \leqslant t < 0, \\ 
\frac{1}{12} & \text{si } 0 \leqslant t < 1, \\ 
\frac{1}{4} & \text{si } 1 \leqslant t < 3, \\ 
\frac{1}{3} & \text{si } 3 \leqslant t < 5, \\ 
1 & \text{si } t \geqslant 5. 
\end{cases}$$

\begin{enumerate}
\item Montrer que $X$ est une v.a. discrète et déterminer sa fonction de masse.
\item Calculer $\mathbb{P}\left(-3 \leqslant X < \frac{1}{2}\right)$.
\end{enumerate}
\end{exercice}
\end{mdframed}

\medskip

\begin{mdframed}
\begin{exercice}
Soit $X$ une variable aléatoire dont la fonction de répartition $F_X$ est définie par :

$$
\forall t \in \mathbb{R} \quad F_X(t) = \begin{cases} 
0 & \text{si } t \leqslant 0, \\ 
\frac{t}{2} & \text{si } 0 \leqslant t \leqslant 2, \\ 
1 & \text{si } t \geqslant 2. 
\end{cases}
$$

\begin{enumerate}
\item Tracer la courbe représentative de $F_X$, montrer que $X$ est une v.a. à densité et déterminer cette densité.
\item Calculer $\mathbb{P}\left(\frac{1}{4} < X < \frac{3}{4}\right)$.
\end{enumerate}
\end{exercice}
\end{mdframed}

\medskip

\begin{mdframed}
\begin{exercice}
Soit $X$ une variable aléatoire de densité $f$ définie par

$$
f(x) = \begin{cases} 
\frac{10}{x^2} & \text{si } x > a, \\ 
0 & \text{si } x \leqslant a, 
\end{cases}
$$

pour un certain $a > 0$.

\begin{enumerate}
\item Calculer $a$.
\item Calculer $\mathbb{P}(X > 20)$.
\item Quelle est la fonction de répartition de $X$ ?
\end{enumerate}
\end{exercice}
\end{mdframed}

\medskip

\begin{mdframed}
\begin{exercice}
Soit $X$ une variable aléatoire de densité $f$ définie par

$$
\forall x \in \mathbb{R} \quad f(x) = \frac{1}{4} \mathbb{1}_{\{0 < x < 1\}} + \frac{3}{8} \mathbb{1}_{\{3 < x < 5\}}.
$$

\begin{enumerate}
\item Donner la fonction de répartition de $X$.
\item Calculer $\mathbb{P}\left(X \in \left[\frac{1}{2}, 2\right] \cup \left]3, 4\right]\right)$.
\end{enumerate}
\end{exercice}
\end{mdframed}